\documentclass{article}
\usepackage[utf8]{inputenc}
\usepackage[spanish]{babel}
\usepackage{listings}
\usepackage{xcolor}

\title{Guía de Desarrollo y Entrega de Módulos - Curso Seguridad Informática}
\author{Coordinación del Proyecto}
\date{\today}

\begin{document}

\maketitle

Esta guía establece los 5 puntos obligatorios para que todos los estudiantes (Grupos 1-4) entreguen sus funciones de auditoría de forma consistente.

\section*{1. Estándar de Docstrings (Estilo Google)}
Cada función debe incluir un comentario inicial que explique su propósito, parámetros y valor de retorno. Esto es vital para el mantenimiento del código.
\begin{lstlisting}[language=Python]
def funcion_ejemplo(ip_objetivo: str):
    """
    Realiza una accion especifica de auditoria.
    Args:
        ip_objetivo (str): La direccion IP a analizar.
    Returns:
        dict: Resultados obtenidos.
    """
\end{lstlisting}

\section*{2. Tipado de Datos (Type Hinting)}
Es obligatorio usar \textit{Type Hints} en Python para evitar errores de tipo. Ejemplo: \texttt{def mi\_funcion(param: str) -> dict:}.

\section*{3. Formato de Salida Estandarizado}
Para facilitar el módulo de reportes de la Semana 10, toda función \textbf{debe retornar un diccionario} (no solo imprimir en consola).
Este debe cumplir con el esquema definido en \texttt{docs/schema\_resultados.json}.
\begin{lstlisting}[language=Python]
return {
    "modulo": "DNS",
    "grupo": 1,
    "estudiante": "E1",
    "target": "google.com",
    "timestamp": "2023-10-27T10:00:00Z",
    "status": "success",
    "data": {"registros_a": ["1.1.1.1"]},
    "error_message": null
}
\end{lstlisting}

\section*{4. Manejo de Excepciones y Errores}
Las funciones no deben detener la ejecución del programa principal. Deben usar bloques \texttt{try-except} para capturar errores de red o permisos y devolver un mensaje de error dentro del diccionario de resultados.

\section*{5. Flujo de Trabajo en Git}
\begin{itemize}
    \item Prohibido modificar \texttt{auditoria.py} directamente.
    \item Trabajar exclusivamente en el archivo asignado a su grupo (ej. \texttt{modulos/scanning.py}).
    \item Antes de entregar, asegúrese de eliminar el \texttt{raise NotImplementedError} de su función.
\end{itemize}

\end{document}